\subsection{定理を使う}
証明が済んだので,定理を実際の極限に使ってみよう.

\begin{bookpoint}
\textbf{POINT}\par
$0$ と $0$,または無限大と無限大の形であることと,分母の導関数が $0$ でないことを確かめてから,分子と分母を別々に微分した商の極限を求める.
\end{bookpoint}

\begin{example}{消える次数を比較する}
$\lim_{x\to0}\cfrac{\log(1+x)-x+\cfrac{x^2}{2}}{x^3}$ を求めよ.
\end{example}

\paragraphbf{方針}
$0$ と $0$ の形である.分子を微分すると $\frac{x^2}{1+x}$ とまとまり,分母の導関数 $3x^2$ と $x^2$ で約分できるので,1回微分した商の極限を求めれば済む.

\paragraphbf{解答}
$\frac13$

\paragraphbf{解法}
$x\to0$ のとき分子も分母も $0$ に近づき,\,$x\ne0$ で分母の導関数 $3x^2$ は $0$ でない.分子を微分すると
\[
 \frac1{1+x}-1+x=\frac{x^2}{1+x}
\]
とまとめられるので,
\[
 \cfrac{\left(\log(1+x)-x+\cfrac{x^2}{2}\right)'}{(x^3)'}
 =\frac{x^2}{3x^2(1+x)}
 =\frac1{3(1+x)}\longrightarrow\frac13
\]
となる.この極限は左右どちらから近づけても $\frac13$ なので,ロピタルの定理により,求める極限は $\frac13$ である.微分した後に式を整理すると $x^2$ が約分できるため,3回も微分する必要はない.

\paragraphbf{条件を省けない例}
ロピタルの定理の条件は省略できない.たとえば,\,$x\to0$ で $\frac{x^2+1}{x^2}$ を考えると,微分した商は $\frac{2x}{2x}=1$ だが,元の商は $+\infty$ へ発散する.元の分子は $1$ に近づき,\,$0$ にも無限大にも近づかないので,定理の仮定を満たしていないのである.

また,定理は逆向きには使えない.\,$x\to+\infty$ で
\[
 \frac{x+\sin x}{x}=1+\frac{\sin x}{x}\longrightarrow1
\]
であるが,微分した商 $1+\cos x$ には極限がない.微分した商の極限が存在しないことから,元の商の極限が存在しないとは結論できない.
